\section{Theoretical Foundations}
\label{sec:theory}

This framework does not invent new ethical categories. Instead, it triangulates a single structural requirement from three independent intellectual traditions. Whether viewed through physics, political economy, or code, the stability of a system depends on the capacity of its subjects to provide corrective feedback.

\subsection{Cybernetics and Requisite Variety}

In cybernetics, the relationship between a controller and its environment is governed by Ashby's Law of Requisite Variety \citep{ashby1956}. The law states that for a system to be stable, the number of control states must meet or exceed the number of environmental states:

\begin{equation}
V(\text{controller}) \geq V(\text{disturbance})
\end{equation}

In DPI, the ``disturbance'' is the boundless diversity of the human population. No pre-programmed rule set can match this variety. Stability, therefore, relies on feedback channels that transmit error signals from the population back into the system's behavior. If these channels (specifically the ability to exit or correct) are blocked, the controller's variety effectively collapses. The system loses the capacity to regulate its own impact, leading inevitably to volatility or oppression.

\subsection{The Commons and Constitutional Constraint}

Elinor Ostrom's work on commons governance \citep{ostrom1990} establishes that sustainable systems require monitors who are accountable to the users and collective-choice arrangements that allow users to modify the rules. A system in which rule-making is the exclusive preserve of the operator acts as an enclosure, not a commons.

Critically, Ostrom identifies that governance is not merely about access; it is about graduated sanctions and low-cost conflict resolution. Applied to digital infrastructure, this implies that users must possess \textbf{Constitutive Power}: the ability to set binding limits on the system's behavior. If a system's constraints are merely advisory or can be overridden by the operator at will, the system is fundamentally ungoverned.

\subsection{Free Software and the Right to Reproduce}

The Free Software movement \citep{stallman2002} contributes the mechanical means of enforcement. Stallman's ``Four Freedoms'' are often misread as a philosophy of sharing, but they are functionally a philosophy of power. The freedom to study (Inspectability) ensures power is visible; the freedoms to modify and distribute (Reproduction) ensure that power is not a monopoly.

This tradition clarifies that openness is not a passive property (reading code) but an active one (forking infrastructure). Transparency without the power to act on it is voyeurism. The ability to reproduce the system independent of its original steward is the only ultimate check against operator capture.